\chapter{Implementation}
\label{chap:implementation}

\section{Implementation trajectory}

Implementation progressed in two explicit stages:

\begin{itemize}
  \item \textbf{Stage 1: n8n phase} requested by the company for rapid feasibility
    validation;
  \item \textbf{Stage 2: new Python platform} for maintainability, stronger memory,
    clearer observability, and reliable long-run execution.
\end{itemize}

The second stage keeps the lessons of the first stage while removing structural
limitations.


\section{Backend and operator surfaces}

The runtime uses FastAPI for orchestration and API exposure, and Next.js for operator
workflows (run launch, live stream of events, historical inspection, and evaluation
review). The operator console streams run events in real time via Server-Sent Events
(SSE), allowing operators to watch each tool call and agent decision as it happens.

\section{Four-agent design: specialist roles and boundaries}

Each agent is designed around one clear responsibility. Figure~\ref{fig:four-agents}
shows the four specialist agents with their tools and interfaces.

\begin{figure}[H]
\centering
\begin{tikzpicture}[node distance=0.4cm and 0.6cm]

% ===== Classification Agent =====
\node[agentbox, minimum height=5cm, minimum width=5.8cm,
      fill=blue!10, draw=blue!60] (cls) at (0, 0) {
};
\node[font=\small\bfseries, text=blue!70] at (0, 1.8) {Classification Agent};
\node[toolbox, minimum width=5cm] at (0, 1.0) {inspect\_page (read-only)};
\node[toolbox, minimum width=5cm] at (0, 0.2) {screenshot()};
\node[toolbox, minimum width=5cm] at (0,-0.6) {get\_frame\_tree()};
\node[toolbox, minimum width=5cm, fill=yellow!10, draw=orange!50] at (0,-1.4)
  {Output: page\_type label};
\node[font=\tiny, text=gray!70] at (0,-2.0) {Decides routing only};

% ===== Landing Agent =====
\node[agentbox, minimum height=5cm, minimum width=5.8cm,
      fill=orange!10, draw=orange!60] (land) at (6.4, 0) {
};
\node[font=\small\bfseries, text=orange!80] at (6.4, 1.8) {Landing Agent};
\node[toolbox, minimum width=5cm] at (6.4, 1.0) {inspect\_landing()};
\node[toolbox, minimum width=5cm] at (6.4, 0.2) {interact (tab click, scroll)};
\node[toolbox, minimum width=5cm] at (6.4,-0.6) {screenshot()};
\node[toolbox, minimum width=5cm, fill=yellow!10, draw=orange!50] at (6.4,-1.4)
  {Output: ranked host URLs};
\node[font=\tiny, text=gray!70] at (6.4,-2.0) {Discovers hosting candidates};

% ===== Hosting Agent =====
\node[agentbox, minimum height=5.6cm, minimum width=5.8cm,
      fill=green!8, draw=green!60] (host) at (0, -6.8) {
};
\node[font=\small\bfseries, text=green!70] at (0, -4.8) {Hosting Agent};
\node[toolbox, minimum width=5cm] at (0,-5.6) {inspect\_hosting()};
\node[toolbox, minimum width=5cm] at (0,-6.4) {interact (server btn, play)};
\node[toolbox, minimum width=5cm] at (0,-7.2) {harvest()};
\node[toolbox, minimum width=5cm] at (0,-8.0) {get\_media\_state()};
\node[toolbox, minimum width=5cm, fill=yellow!10, draw=orange!50] at (0,-8.8)
  {Output: m3u8/mpd streams};
\node[font=\tiny, text=gray!70] at (0,-9.4) {Activates players, extracts streams};

% ===== Embedded Agent =====
\node[agentbox, minimum height=5.6cm, minimum width=5.8cm,
      fill=purple!8, draw=purple!60] (emb) at (6.4, -6.8) {
};
\node[font=\small\bfseries, text=purple!70] at (6.4, -4.8) {Embedded Agent};
\node[toolbox, minimum width=5cm] at (6.4,-5.6) {inspect\_embedded()};
\node[toolbox, minimum width=5cm] at (6.4,-6.4) {interact (play btn, frame)};
\node[toolbox, minimum width=5cm] at (6.4,-7.2) {harvest()};
\node[toolbox, minimum width=5cm] at (6.4,-8.0) {navigate()};
\node[toolbox, minimum width=5cm, fill=yellow!10, draw=orange!50] at (6.4,-8.8)
  {Output: final stream manifests};
\node[font=\tiny, text=gray!70] at (6.4,-9.4) {Fallback: iframe player contexts};

% Flow arrows
\draw[myarrow, very thick] (cls.east) -- node[above, font=\tiny]{host\_page}
  (land.west);
\draw[myarrow, very thick] (cls.south) -- node[right, font=\tiny]{host}
  (host.north);
\draw[myarrow, very thick] (land.south) -- node[right, font=\tiny]{host URLs}
  (emb.north);
\draw[->, >=Stealth, thick, red!60, dashed] (host.east) --
  node[above, font=\tiny, red!60]{fallback} (emb.west);

\end{tikzpicture}
\caption{Four specialist agents: tools, outputs, and routing relationships.}
\label{fig:four-agents}
\end{figure}

\section{Agent design: reasoning toward a final objective}

Each agent is built around a single concrete objective rather than a general instruction to "do useful things." The classification agent needs to produce a routing label — nothing else. The landing agent needs to find all hosting-page candidates on a listing site, including ones behind tabs and scrolled content. The hosting agent's job is to activate server controls, deal with whatever popups appear, and return stream URLs; if it can't find a stream after exhausting its options, it returns the embed URL so the pipeline can try the embedded agent. The embedded agent handles the fallback case: iframe-origin player pages where the main context had nothing to offer.

\section{Prompt design and structure per agent}

This section documents the prompt structure for each agent. Prompts evolved from
intent-based to ReAct-style over the course of the internship. The structures shown
here reflect the final ReAct-style design.

\subsection{Classification agent prompt}

The classification agent receives a minimal toolset and must produce a confident
page-type label with brief reasoning. It is the cheapest agent to run and uses the
lightweight model.

The prompt has three sections. First, a \emph{role and objective} block that constrains the agent to a
single classification output — one of \texttt{landing\_page}, \texttt{host\_page},
\texttt{embed\_video\_page}, or \texttt{other} — returned as a JSON object with a confidence
score and reasoning field. Second, a \emph{procedure} block that enforces a strict two-tool
sequence: call \texttt{inspect\_page} for structural DOM signals, then \texttt{screenshot}
for visual confirmation, then reason and return. Third, an \emph{important} block that caps
tool usage at three calls and prohibits any page interaction, keeping classification fast
and inexpensive.

\subsection{Landing-page agent prompt}

The landing agent must navigate listing pages to discover all hosting-page candidates.
Its key challenge is that some matches only appear after interacting with date tabs or
scrolling.

The prompt encodes a full
ReAct loop: \textsc{observe} with \texttt{inspect\_landing}, \textsc{think} about which
links are hosting candidates and whether any unexplored date tabs remain, \textsc{act} by
clicking tabs or scrolling, and \textsc{reflect} on completeness before repeating. Stopping
criteria are explicit: all visible tabs inspected, eight consecutive tool calls without new
candidates, or twenty candidates already found. This prevents runaway exploration on
paginated listing sites.

\subsection{Hosting-page agent prompt}

The hosting agent is the most complex. It must activate server controls, handle popups,
validate player state after each interaction, and fall back to the embedded agent when
needed.

The prompt is the longest of the
four. Its ReAct loop begins with \texttt{inspect\_hosting} to map server buttons, player
containers, and any pre-loaded network signals. If no \texttt{.m3u8} pattern is already
visible in the network, the agent activates the primary server button via \texttt{interact},
then calls \texttt{screenshot} to confirm the click effect before invoking \texttt{harvest}.
A dedicated popup-handling block instructs the agent to detect new tabs or overlays and
retry with an alternative click strategy (\texttt{text\_content} or \texttt{js\_click})
rather than failing. If all server options are exhausted without a stream, the agent returns
the embedded player URL so the orchestrator can route to the embedded agent. The budget cap
is 12 tool calls, reflecting the higher complexity of this agent's task.

\subsection{Embedded-page agent prompt}

The embedded agent works within iframe-origin player pages. Its main task is to
trigger play and capture network-level stream evidence.

The prompt reflects the
narrower scope of an embedded player page: no navigation, no server cycling, just player
activation and stream capture. The ReAct loop checks whether the player is auto-playing
(call \texttt{harvest} immediately) or requires a play click (dismiss any ad overlay first,
then click, wait two seconds for network requests to materialize, then \texttt{harvest}).
A detailed player activation pattern section covers selector-based, coordinate-based, and
JavaScript-injected click strategies in priority order. The budget is capped at three
interaction attempts before the agent declares failure and returns a \texttt{stream\_masked}
or \texttt{no\_play\_button} failure reason for the operator to review.

\section{Puppeteer tool registry and tool descriptions}

\subsection{Why tool descriptions matter}

When an LLM agent decides which tool to call next, it reads each tool's schema
description. This is not cosmetic documentation --- it is the primary signal the
model uses to distinguish between similar-sounding tools. A vague description
causes wrong tool selection; a precise description causes the model to call the
right tool on the first attempt.

This behavior is also API-provider specific. Some providers (Google Gemini~\cite{gemini},
Anthropic Claude~\cite{claude}) weight tool descriptions heavily in tool-selection
decisions. A well-phrased description can eliminate entire categories of tool-call
mistakes without any prompt change.

\textbf{Design rule applied in this project:} every tool description must answer three
questions: \emph{what state it reads or modifies}, \emph{what signals it returns},
and \emph{when to use it versus a similar tool}.

\subsection{Tool registry overview}

The Puppeteer tool layer exposes the following tools to all agents, with descriptions
that shape agent behavior:

\begin{table}[H]
\centering
\begin{tabularx}{\textwidth}{>{\raggedright\arraybackslash\bfseries}p{3cm}
                              >{\raggedright\arraybackslash}X}
\toprule
\rowcolor{accent!8}
Tool name & Description exposed to the LLM \\
\midrule
\texttt{inspect\_page} &
  Read the current page DOM state. Returns title, meta description, visible text links,
  iframe src list, player signals (video/source tags, player library class names), and
  network request summary. Use this as the first action after any navigation. Do not
  use to activate controls. \\
\texttt{inspect\_landing} &
  Specialized inspect for listing pages. Returns match link candidates, date/category
  tab labels, broadcaster references, and pagination controls. Use instead of
  inspect\_page when the classification result is landing\_page. \\
\texttt{inspect\_hosting} &
  Specialized inspect for host pages. Returns server button selectors and labels,
  player container state, iframe src attributes, and current network stream signals.
  Use instead of inspect\_page when on a host page. \\
\texttt{inspect\_embedded} &
  Specialized inspect for embedded player pages. Returns player library type
  (albaplayer, clappr, hls.js, videojs), play button presence and selector, auto-play
  status, and frame nesting context. Use when operating within an iframe-origin page. \\
\texttt{interact} &
  Execute a single user interaction on the page. Strategies: selector, text\_content,
  xpath, coordinates, js\_click. Returns action outcome and whether page state changed.
  Always follow with screenshot() to confirm. Never call harvest() before interact()
  if the player is not yet playing. \\
\texttt{screenshot} &
  Capture a PNG screenshot of the current page and upload to Cloudinary. Returns the
  CDN URL. Use after every interact() call and before every harvest() call to confirm
  state. \\
\texttt{harvest} &
  Collect stream evidence from all available layers: CDP network event log, JavaScript
  player source attributes, and performance.getEntriesByType retroactive scan. Returns
  m3u8/mpd/mp4 URL candidates with source layer attribution. Use only after player
  activation is confirmed by screenshot or get\_media\_state. \\
\texttt{get\_media\_state} &
  Probe the current player state by evaluating JavaScript player API objects in the
  active frame. Returns: playing, paused, idle, error, or unknown. Use before harvest()
  to decide whether activation is needed. \\
\texttt{navigate} &
  Navigate the browser to a new URL with a configurable wait strategy (load,
  networkidle0, networkidle2). Returns the final URL after redirects. Use to follow
  embed URLs or explicit server switches that require a full page load. \\
\texttt{get\_frame\_tree} &
  Enumerate all frames in the current page, including nested iframes, with their origin
  URLs and nesting depth. Use when inspect\_page shows iframe signals but the player
  frame is not accessible from the top context. \\
\bottomrule
\end{tabularx}
\caption{Puppeteer tool descriptions as exposed to the LLM agent.}
\label{tab:tool-descriptions}
\end{table}

\subsection{How tool descriptions guide agent behavior: concrete examples}

\textbf{Example 1: inspect\_page vs inspect\_hosting.}
Before the specialized inspect tools were added, agents called \texttt{inspect\_page}
on host pages and received generic DOM information that did not highlight server
buttons. Adding \texttt{inspect\_hosting} with the description ``Returns server button
selectors and labels'' caused the agent to switch tools automatically on host pages
without any prompt instruction to do so.

\textbf{Example 2: harvest timing.}
The \texttt{harvest} description includes the phrase ``Use only after player activation
is confirmed.'' This sentence alone reduced premature harvest calls (calling harvest
before clicking play) by a large margin compared to the early prompt versions that
relied on the system prompt instruction alone.

\textbf{Example 3: get\_media\_state disambiguation.}
Before \texttt{get\_media\_state} existed, agents used \texttt{screenshot} as their
only state-check tool. This worked visually but was slower. Adding
\texttt{get\_media\_state} with the description ``Probe the current player state by
evaluating JavaScript player API objects'' caused agents to use it as a fast pre-check
before the more expensive screenshot call.

\subsection{Per-agent tool allocation}

Not all tools are available to all agents. Each agent receives only the tools relevant
to its role, which reduces the decision space and prevents cross-agent tool misuse.

\begin{table}[H]
\centering
\begin{tabularx}{\textwidth}{>{\raggedright\arraybackslash\bfseries}p{3.8cm}
                              >{\raggedright\arraybackslash}X}
\toprule
\rowcolor{accent!8}
Agent & Tools available \\
\midrule
Classification agent &
  \texttt{inspect\_page}, \texttt{screenshot}, \texttt{get\_frame\_tree} \\
Landing agent &
  \texttt{inspect\_landing}, \texttt{interact}, \texttt{screenshot} \\
Hosting agent &
  \texttt{inspect\_hosting}, \texttt{interact}, \texttt{screenshot},
  \texttt{get\_media\_state}, \texttt{harvest}, \texttt{navigate},
  \texttt{get\_frame\_tree} \\
Embedded agent &
  \texttt{inspect\_embedded}, \texttt{interact}, \texttt{screenshot},
  \texttt{harvest}, \texttt{navigate} \\
\bottomrule
\end{tabularx}
\caption{Tool allocation per specialist agent. Agents receive only tools relevant to
their role to reduce tool-selection error surface.}
\label{tab:per-agent-tools}
\end{table}

\subsection{Diagram: tool call flow per agent type}

Figure~\ref{fig:tool-flow} shows the typical tool call sequence for each agent,
reflecting the order in which tools are described in the prompt procedure.

\begin{figure}[H]
\centering
\begin{tikzpicture}[
  seq/.style={draw=accent!70, fill=blue!5, rounded corners=3pt,
              minimum width=2.9cm, minimum height=0.7cm,
              font=\tiny, align=center},
  arr/.style={->, >=Stealth, thin, gray!70}
]

% Classification
\node[font=\footnotesize\bfseries, text=blue!70] at (1.5, 5.2) {Classification};
\node[seq] (c1) at (1.5, 4.6) {inspect\_page};
\node[seq] (c2) at (1.5, 3.7) {screenshot};
\node[seq, fill=green!8, draw=green!60] (c3) at (1.5, 2.8) {[return type]};
\draw[arr](c1)--(c2); \draw[arr](c2)--(c3);

% Landing
\node[font=\footnotesize\bfseries, text=orange!80] at (5.2, 5.2) {Landing};
\node[seq] (l1) at (5.2, 4.6) {inspect\_landing};
\node[seq] (l2) at (5.2, 3.7) {interact (tab)};
\node[seq] (l3) at (5.2, 2.8) {inspect\_landing};
\node[seq, fill=green!8, draw=green!60] (l4) at (5.2, 1.9) {[return URLs]};
\draw[arr](l1)--(l2); \draw[arr](l2)--(l3); \draw[arr](l3)--(l4);

% Hosting
\node[font=\footnotesize\bfseries, text=green!70] at (8.9, 5.2) {Hosting};
\node[seq] (h1) at (8.9, 4.6) {inspect\_hosting};
\node[seq] (h2) at (8.9, 3.7) {interact (server)};
\node[seq] (h3) at (8.9, 2.8) {screenshot};
\node[seq] (h4) at (8.9, 1.9) {get\_media\_state};
\node[seq] (h5) at (8.9, 1.0) {harvest};
\node[seq, fill=green!8, draw=green!60] (h6) at (8.9, 0.1) {[return streams]};
\draw[arr](h1)--(h2); \draw[arr](h2)--(h3);
\draw[arr](h3)--(h4); \draw[arr](h4)--(h5); \draw[arr](h5)--(h6);

% Embedded
\node[font=\footnotesize\bfseries, text=purple!70] at (12.6, 5.2) {Embedded};
\node[seq] (e1) at (12.6, 4.6) {inspect\_embedded};
\node[seq] (e2) at (12.6, 3.7) {interact (play)};
\node[seq] (e3) at (12.6, 2.8) {screenshot};
\node[seq] (e4) at (12.6, 1.9) {harvest};
\node[seq, fill=green!8, draw=green!60] (e5) at (12.6, 1.0) {[return streams]};
\draw[arr](e1)--(e2); \draw[arr](e2)--(e3);
\draw[arr](e3)--(e4); \draw[arr](e4)--(e5);

% Separator lines
\foreach \x in {3.35, 7.05, 10.75}
  \draw[dashed, gray!40] (\x, 5.5) -- (\x, -0.3);

\end{tikzpicture}
\caption{Typical tool call sequence per agent. Each column shows the ordered tool
invocations; agents loop back when the stopping criterion is not yet met.}
\label{fig:tool-flow}
\end{figure}

\section{Tool architecture and why it works on hard sites}

The tool architecture combines high-level operations (inspect/interact/harvest/
screenshot/navigate) with lower-level primitives (frame tree, page context, media state,
action variants). This layered design helps agents adapt to unknown page structures
instead of depending on one brittle selector path.

\begin{table}[H]
\centering
\begin{tabularx}{\textwidth}{>{\raggedright\arraybackslash\bfseries}p{3.2cm}
                              >{\raggedright\arraybackslash}X
                              >{\raggedright\arraybackslash}X}
\toprule
\rowcolor{accent!8}
Observed obstacle & Typical site behavior & Tool-layer adaptation \\
\midrule
Theme-level selector drift & same visual page, changed underlying selectors &
  text and role-based fallback targeting, not selector-only targeting \\
Mirror heterogeneity & one event points to multiple providers with different UI
  semantics & provider-agnostic inspect/interact loops with per-mirror diagnostics \\
Ad diversion & interaction causes tab hijack or overlay interruption &
  focus recovery, popup handling, and retry with state validation \\
Late stream materialization & stream appears only after runtime player logic &
  deferred harvest triggered by confirmed media-state transition \\
\bottomrule
\end{tabularx}
\caption{How the tool architecture adapts to recurrent break behaviors.}
\label{tab:tool-adaptation}
\end{table}

\section{Puppeteer techniques inside tools}

Key techniques used in tools include:

\begin{itemize}
  \item \texttt{page.evaluate(...)} for runtime extraction of player state and
    DOM-derived signals;
  \item CDP network events for request/response-level evidence capture;
  \item iframe diagnostics and frame-targeted interactions;
  \item performance resource scanning for retroactive stream evidence;
  \item screenshot checkpoints after each critical state transition.
\end{itemize}

These methods are essential because static HTML alone does not expose enough evidence
in many real cases.

\section{Handling tokenized \texttt{m3u8} and hidden stream flows}

HLS (HTTP Live Streaming)~\cite{hls-rfc} manifests (\texttt{.m3u8}) are the primary
evidence target in this project. Some websites do not expose final manifests directly:
the stream appears only after runtime token exchange or player-side JavaScript logic. In such runs, an initial page
resource can look unrelated (including decoy assets like \texttt{mono.css}) while the
real signed stream URL is generated later.

Implementation strategy:

\begin{itemize}
  \item trigger real player behavior first (interaction step);
  \item capture network request/response layers after activation;
  \item inspect JavaScript player objects and runtime resource entries;
  \item aggregate and deduplicate candidate stream URLs;
  \item validate with screenshots and media-state diagnostics.
\end{itemize}

\section{Screenshot capture and Cloudinary storage}

Screenshot reasoning was treated as a first-class mechanism, not a cosmetic artifact.
It serves three purposes:

\begin{itemize}
  \item confirms whether an interaction changed state (player started, popup closed,
    server switched);
  \item prevents false-positive claims when a click did nothing visible;
  \item provides operator-auditable evidence for manual verification and takedown
    package assembly.
\end{itemize}

\textbf{Cloudinary integration.} After each \texttt{screenshot()} call, the Puppeteer
tool uploads the PNG buffer to Cloudinary using their unsigned upload preset. The tool
returns the Cloudinary CDN URL to the agent, which stores it as part of the run event
record. The URL is then surfaced in the operator console next to the corresponding
tool call in the live trace.

This approach decouples evidence storage from the extraction container. Containers can
be ephemeral (as required by Azure Container Apps Jobs) while screenshots remain
accessible indefinitely via the Cloudinary CDN link.

\section{Model strategy, cache behavior, and cost}

The implementation uses smaller and larger models with role-based allocation:

\begin{itemize}
  \item \texttt{gemini-2.5-flash-lite} for orchestration and lightweight routing
    decisions;
  \item \texttt{gemini-2.5-flash} for specialist extraction where richer reasoning is
    required.
\end{itemize}

Observed cost behavior is strongly model-dependent and task-depth dependent.

\begin{table}[H]
\centering
\begin{tabularx}{\textwidth}{>{\raggedright\arraybackslash\bfseries}p{3.7cm}
                              >{\raggedright\arraybackslash}X
                              >{\raggedright\arraybackslash}X}
\toprule
\rowcolor{accent!8}
Cost driver & Why it increases cost & Mitigation used \\
\midrule
Model choice & stronger models cost more per token &
  route lightweight tasks to flash-lite \\
Hosting-page count & more pages require more reasoning/tool loops &
  ranking and early pruning \\
Servers per page & each server adds interaction + harvest attempts &
  bounded retries and evidence thresholds \\
Long context windows & repeated prompt context can become expensive &
  provider cache and tool-result cache \\
\bottomrule
\end{tabularx}
\caption{Primary cost drivers and implementation mitigations.}
\label{tab:cost-drivers}
\end{table}

\section{Implementation summary}

The real test of the implementation was not in isolation but under actual streaming-site conditions: nested iframes, popup click traps, tokenized manifests that don't appear until 2+ seconds after play activation, and Cloudflare challenges that serve alternate HTML to anything that looks like a bot. The layered tool architecture handled most of these without prompt changes; the explicit fallback from hosting to embedded handled the rest.

The architecture is also compatible with queue-based cloud execution through Azure Service Bus and Container Apps Jobs. That deployment model was designed and documented during the internship but daily execution ran locally and in containers, which proved fast enough for evaluation campaigns.

Chapter~6 evaluates these design decisions through concrete use cases and prompt evolution measurements, including the impact each prompt version had on navigation recall, player activation rate, and overall stream recovery.
